\title{DeepSeekMath: Pushing the Limits of Mathematical Reasoning in Open Language Models}

\begin{abstract}
Mathematical reasoning remains a substantial hurdle for language models, primarily due to its inherent complexity and structured logic. In this study, we present DeepSeekMath~7B, an extension of DeepSeek-Coder-Base-v1.5 7B, further pre-trained on 120B mathematics-focused tokens collected from Common Crawl, as well as additional natural language and code resources. DeepSeekMath~7B demonstrates remarkable performance by attaining a 51.7\% score on the challenging, competition-grade MATH benchmark—accomplished without the aid of external tools or ensemble voting—thus nearing the proficiency levels observed in Gemini-Ultra and GPT-4. Furthermore, employing self-consistency across 64 samples, DeepSeekMath~7B achieves a 60.9\% accuracy rate on MATH. Two principal innovations underlie DeepSeekMath’s robust mathematical reasoning: First, we exploit large-scale, openly available web corpora through a carefully crafted data selection strategy. Second, we propose Group Relative Policy Optimization (GRPO), a novel adaptation of Proximal Policy Optimization (PPO), which boosts the model’s mathematical reasoning competence while simultaneously optimizing PPO’s memory consumption.
\end{abstract}

\section{Introduction}

The advent of large language models (LLMs) has transformed methods for tackling mathematical reasoning in artificial intelligence, yielding notable improvements on quantitative benchmarks \citep{MATH} and geometric reasoning tasks \citep{trinh2024solving}. In addition, these LLMs have become valuable tools for human users grappling with intricate mathematical challenges \citep{tao}. Nonetheless, leading-edge systems like GPT-4 \citep{gpt4} and Gemini-Ultra \citep{gemini} remain closed-source, leaving available open-source models lagging significantly behind in terms of performance.

In this paper, we present DeepSeekMath, a specialized language model tailored for mathematical tasks that not only surpasses open-source alternatives in mathematical reasoning but also nearly matches the performance of GPT-4 on established academic benchmarks. Central to our approach is the construction of the DeepSeekMath Corpus, a comprehensive and high-quality pre-training dataset containing 120B mathematical tokens. We source this data from Common Crawl (CC) utilizing a classifier based on fastText \citep{joulin2016fasttext}. Initially, this classifier is trained with positive samples drawn from OpenWebMath \citep{openwebmath} and a broad assortment of unrelated web pages as negative samples. In subsequent steps, we deploy the trained classifier to identify further mathematical content from CC, which undergoes additional refinement through manual annotation. The enhanced annotations are then incorporated to retrain and improve the classifier. Our evaluation demonstrates that the resulting corpus is of substantial quality: the DeepSeekMath-Base 7B model achieves 64.2\% accuracy on GSM8K \citep{gsm8k} and 36.2\% on the MATH competition benchmark \citep{MATH}, exceeding the results of Minerva 540B \citep{minerva}. Furthermore, since the DeepSeekMath Corpus covers multiple languages, we also observe performance gains on Chinese mathematical datasets \citep{wei2023cmath,agieval}. We suggest that our strategies for mathematical data collection and processing may serve as a valuable foundation for future research, highlighting ample opportunities for further advancement.

DeepSeekMath-Base is built upon the DeepSeek-Coder-Base-v1.5 7B model \citep{deepseek-coder}, given our observation that commencing with a model pre-trained on code data yields superior results compared to initializing with a standard large language model. Additionally, we note that training on mathematical data not only enhances performance on mathematical tasks but also leads to gains on general benchmarks such as MMLU \citep{mmlu} and BBH \citep{bbh}, suggesting that the improvements extend to broader reasoning skills.

Subsequent to pre-training, we fine-tune DeepSeekMath-Base for mathematical reasoning using instruction data encompassing chain-of-thought reasoning \citep{cot}, program-of-thought approaches \citep{pot,pal}, and tool-assisted reasoning \citep{tora}. The resultant DeepSeekMath-Instruct 7B model surpasses all existing 7B-scale models and approaches the performance of instruction-tuned open-source models at the 70B parameter scale.

We further present Group Relative Policy Optimization (GRPO), a novel variant of reinforcement learning based on Proximal Policy Optimization (PPO) \citep{schulman2017proximal}. Unlike conventional PPO, GRPO eliminates the need for a separate critic network and instead utilizes group-based scoring to estimate the baseline, thereby reducing computational overhead. Employing GRPO with only a portion of the English instruction tuning dataset, we observe marked improvements over the robust DeepSeekMath-Instruct model. These include significant performance boosts in both in-domain mathematical evaluations (GSM8K: 82.9\% $\rightarrow$ 88.2\%, MATH: 46.8\% $\rightarrow$ 51.7\%) and out-of-domain assessments (e.g., CMATH: 84.6\% $\rightarrow$ 88.8\%) during the RL phase. 

Furthermore, we propose a unified analytical framework that encompasses methods such as Rejection Sampling Fine-Tuning (RFT) \citep{yuan2023scaling}, Direct Preference Optimization (DPO) \citep{dpo}, PPO, and GRPO, interpreting them as direct or streamlined forms of reinforcement learning. Within this framework, we systematically explore various aspects, including comparisons between online and offline RL training, outcome-based versus process-based supervision, and the impact of single-turn versus iterative RL methods, to identify key determinants underlying the efficacy of each approach. Finally, we articulate how our RL methodology enhances instruction-tuned model performance and discuss future directions for advancing reinforcement learning under this comprehensive paradigm.

\subsection{Contributions}
This work presents advances in large-scale mathematical pre-training and investigates the impacts of reinforcement learning.

\textbf{Math Pre-Training at Scale}

\begin{itemize}[topsep=0pt]
    \item We demonstrate that the openly available Common Crawl data harbors substantial resources for mathematical research. Through a rigorously crafted data selection pipeline, we curate the DeepSeekMath Corpus—a comprehensive, high-fidelity dataset encompassing 120B tokens from mathematics-oriented web pages. This corpus is nearly seven times larger than the math web dataset utilized by Minerva \citep{minerva} and surpasses the recently introduced OpenWebMath \citep{openwebmath} by a factor of nine.

    \item The DeepSeekMath-Base 7B model, pre-trained on this corpus, matches the performance level of the substantially larger Minerva 540B model \citep{minerva}. This outcome reveals that sheer model size is not the sole determinant of mathematical reasoning proficiency; smaller models trained on carefully filtered, high-quality data can demonstrate competitive results.

    \item We report observations from our experimental suite on math pre-training. Preceding mathematical training with code-related training demonstrably enhances a model’s capacity to tackle math problems, irrespective of tool assistance. This provides empirical support for the hypothesis that code pre-training bolsters reasoning skills, at least within the mathematical domain.

    \item Contrary to prevalent practice—particularly the inclusion of arXiv material in many mathematics-related datasets—we find that pre-training on arXiv papers does not yield measurable gains on any of the math evaluation benchmarks considered in this study.
\end{itemize}

\textbf{Exploration and Analysis of Reinforcement Learning}

\begin{itemize}[topsep=0pt]
    \item This work presents Group Relative Policy Optimization (GRPO), a reinforcement learning algorithm that balances effectiveness with efficiency. Unlike conventional methods that depend on a critic model, GRPO leverages group score-derived baselines, resulting in a substantial reduction in computational and memory requirements relative to Proximal Policy Optimization (PPO).
    \item Our empirical studies reveal that incorporating GRPO markedly boosts the capabilities of DeepSeekMath-Instruct—our instruction-tuned language model—when trained solely on instruction-tuning data. Additionally, we document improvements in handling tasks outside the training distribution throughout the reinforcement learning phase.
    \item We propose a unified framework for interpreting a range of methods, including RFT, DPO, PPO, and GRPO. A thorough experimental investigation covers aspects such as online versus offline learning, outcome-based versus process-based supervision, and comparisons between single-step and iterative reinforcement learning. These experiments illuminate the core factors underlying the effectiveness of this unified paradigm.
    \item Building on the proposed framework, we investigate underlying mechanisms that make reinforcement learning effective for large language models, and outline several promising research avenues for enhancing LLM reinforcement learning techniques.
\end{itemize}

\subsection{Summary of Evaluations and Metrics}

\begin{itemize}[topsep=0pt]
    \item \textbf{English and Chinese Mathematical Reasoning}: To evaluate our models, we conduct large-scale testing on a suite of English and Chinese mathematics benchmarks, covering questions from primary to tertiary levels. English evaluation sets comprise GSM8K \citep{gsm8k}, MATH \citep{MATH}, SAT \citep{llemma}, OCW Courses \citep{minerva}, and MMLU-STEM \citep{mmlu}. For Chinese, we use MGSM-zh \citep{mgsm}, CMATH \citep{wei2023cmath}, Gaokao-MathCloze \citep{agieval}, and Gaokao-MathQA \citep{agieval}. Performance is assessed in terms of both the ability to generate fully self-contained, tool-free written solutions, as well as in solving problems with Python.

    On the English datasets, DeepSeekMath-Base demonstrates comparable results to the proprietary Minerva 540B \citep{minerva}, and consistently outperforms all available open-source models such as Mistral 7B \citep{mistral} and Llemma-34B \citep{llemma}, independent of their mathematical pre-training status, often by a considerable gap. Remarkably, DeepSeekMath-Base achieves even better results on the Chinese tasks, which can be attributed to the inclusion of high-quality non-English data and the avoidance of the English-only pre-training practice common in previous efforts \citep{minerva,llemma}. With further instruction tuning and reinforcement learning, DeepSeekMath-Instruct and DeepSeekMath-RL achieve robust gains, surpassing the 50\% accuracy mark on the challenging competition-level MATH dataset—an open-source first.

\item \textbf{Formal Mathematics}: DeepSeekMath-Base is assessed using the informal-to-formal theorem proving benchmark from \citep{dsp_proof}, applied to the miniF2F suite \citep{minif2f} with Isabelle \citep{isabelle} serving as the proof assistant. The model exhibits robust few-shot autoformalization capabilities.
\item \textbf{Natural Language Understanding, Reasoning, and Code}: To provide a holistic view of the model’s proficiency in comprehension, reasoning, and programming, we assess DeepSeekMath-Base on several benchmarks. The Massive Multitask Language Understanding (MMLU) evaluation \citep{mmlu} features 57 multiple-choice tasks spanning a wide array of disciplines. BIG-Bench Hard (BBH) \citep{bbh} introduces 23 intricate tasks, primarily designed to assess multi-step reasoning. Additionally, HumanEval \citep{codex} and MBPP \citep{mbpp} are employed to gauge the model's coding abilities. Pre-training on mathematical data improves the model’s outcomes across both language understanding and reasoning benchmarks.

\section{Math Pre-Training}

\subsection{Data Collection and Decontamination}
This section details the methodology employed to assemble the DeepSeekMath Corpus from Common Crawl data. Figure \ref{fig:data_collect} illustrates our iterative pipeline, which methodically expands a large-scale mathematical dataset from Common Crawl, beginning with a seed dataset comprised of high-quality mathematical material. Notably, this methodology is adaptable for other specialized fields, such as computer programming.

Initially, we utilize OpenWebMath \citep{openwebmath}—a curated selection of high-quality mathematical texts found online—as our starting seed corpus. To broaden this base, we train a fastText classifier \citep{joulin2016fasttext} aimed at retrieving additional web pages resembling OpenWebMath in content. For training, we draw 500,000 samples at random from the seed as positive cases, and pair them with 500,000 negative examples from Common Crawl. We rely on the official open-source fastText library\footnote{\url{https://fasttext.cc}} for implementation, setting parameters as follows: vector dimension at 256, learning rate of 0.1, a maximum n-gram length of 3, a minimum word occurrence threshold of 3, and 3 training epochs. To manage the sheer scale of Common Crawl, we apply URL-level deduplication along with near-duplicate filtering, which distills the dataset to approximately 40B HTML web pages. Employing the trained fastText model, we subsequently identify mathematical web content within this filtered collection. To ensure quality, the web pages are ranked using their model-predicted scores, and only the highest-scoring subset is retained. Pre-training trials using datasets of 40B, 80B, 120B, and 160B top-ranked tokens inform the optimal data volume to keep. For the first cycle, we select the leading 40B tokens for further use.

Following the initial round of data collection, a significant portion of mathematical web pages remained uncaptured, primarily due to the limited diversity present in the positive training examples used for the fastText model. To address this issue, we expanded the seed dataset by incorporating more mathematical web sources, aiming to improve the performance of fastText. Our methodology involved segmenting the entire Common Crawl dataset into mutually exclusive domains, where each domain encompasses web pages sharing a common base URL. For each domain, we determined the proportion of web pages retrieved during the first iteration. If a domain exhibited over 10\% retrieval—indicating mathematical relevance (e.g., \url{mathoverflow.net})—we proceeded to further processing. We then conducted manual annotation of URLs within these math-related domains to identify those associated with mathematical content (such as \url{mathoverflow.net/questions}). Pages linked to these annotated URLs but absent from the initial collection were appended to the seed corpus. Through this iterative process, the number of positive samples increases, enabling fastText to be retrained for higher recall of mathematical data in subsequent rounds. By the fourth iteration, we amassed 35.5M mathematical web pages, corresponding to a cumulative total of 120B tokens. Observing that nearly 98\% of the data collected in the fourth cycle overlapped with that of the third, we concluded that additional iterations would yield diminishing returns and thus ended the collection phase.

To prevent contamination of evaluation benchmarks, we adopt the filtering methodology of \cite{deepseek-coder}, excluding from our dataset any web pages that contain questions or answers derived from English mathematical benchmarks such as GSM8K \citep{gsm8k}, MATH \citep{MATH}, as well as Chinese datasets including CMATH \citep{wei2023cmath} and AGIEval \citep{agieval}. Specifically, any segment of text matching exactly with a 10-gram substring from the benchmarks is excised from the training corpus. For benchmark passages shorter than 10 grams but at least 3 grams, we apply a strict exact-match filtering strategy to further eliminate contaminated data.

\subsection{Validating the Quality of the DeepSeekMath~Corpus}

We conducted a series of pre-training studies to evaluate the DeepSeekMath~Corpus in comparison to several state-of-the-art mathematical text corpora released recently:

\begin{itemize}[topsep=0pt]
    \item \textbf{MathPile} \citep{mathpile}: A corpus composed of multiple sources (8.9B tokens), including textbooks, Wikipedia, ProofWiki, CommonCrawl, StackExchange, and arXiv—with arXiv comprising over 85\% of the content;
    \item \textbf{OpenWebMath} \citep{openwebmath}: This dataset consists of mathematical web content extracted from CommonCrawl, amounting to 13.6B tokens;
    \item \textbf{Proof-Pile-2} \citep{llemma}: An aggregation of mathematical material, merging OpenWebMath, AlgebraicStack (which contributes 10.3B tokens of mathematical code), and arXiv papers (accounting for 28.0B tokens). In our experiments utilizing Proof-Pile-2, we align with \cite{llemma} by adhering to the specified ratio of arXiv:Web:Code as 2:4:1.
\end{itemize}

\subsubsection{Training Setting}
\label{sec:quality-policy}
We adapt a general-purpose language model comprising 1.3B parameters—utilizing the same architecture as the DeepSeek LLMs \citep{deepseek-llm}, hereafter referred to as DeepSeek-LLM 1.3B—by conducting targeted math-domain training. Each model is independently fine-tuned on a respective mathematical dataset for a total of 150B tokens. Training takes place within the resource-efficient HAI-LLM \citep{haillm} platform. Consistent with the protocols established for DeepSeek LLMs, optimization employs AdamW \citep{adamW} configured with $\beta_1=0.9$, $\beta_2=0.95$, and $\mathrm{weight\_decay}=0.1$. The learning rate follows a staged decay: it ramps up to its maximum over 2,000 warmup steps, reduces to 31.6\% of the maximum by the 80\% mark of training, and further drops to 10.0\% at the 90\% milestone. The upper bound for the learning rate is set to 5.3e-4. Mini-batches contain 4M tokens, with each sample having a context length capped at 4K tokens.

\subsubsection{Evaluation Results}

\textbf{The DeepSeekMath~Corpus distinguishes itself by its exceptional quality, extensive multilingual content, and its unprecedented scale among available mathematical corpora.}

\begin{itemize}[topsep=0pt]
    \item \textbf{High-quality}:
    Model performance is assessed across 8 mathematical benchmarks utilizing few-shot chain-of-thought prompting \cite{cot}. As detailed in Table \ref{tab:corpora_comparison}, models trained on DeepSeekMath~Corpus consistently outperform alternatives. Figure \ref{fig:corpora_comparisons} further demonstrates that a model fine-tuned with the DeepSeekMath~Corpus achieves superior results over one trained on Proof-Pile-2 after 50B tokens (equivalent to one complete epoch on Proof-Pile-2), highlighting the superior average quality of the DeepSeekMath~Corpus.
    \item \textbf{Multilingual}:
    DeepSeekMath~Corpus comprises data in various languages, with English and Chinese making up the largest portions. According to Table \ref{tab:corpora_comparison}, incorporating DeepSeekMath~Corpus during training leads to enhanced mathematical reasoning capabilities in both English and Chinese. Conversely, most existing math corpora are predominantly English-based and confer minimal, or even negative, impact on reasoning performance for Chinese content.
    \item \textbf{Large-scale}:
    In terms of size, DeepSeekMath~Corpus surpasses current math-focused corpora by a considerable margin. Figure \ref{fig:corpora_comparisons} illustrates that DeepSeek-LLM 1.3B, trained on DeepSeekMath~Corpus, exhibits a steeper and more sustained learning trajectory. Baseline datasets, on the other hand, are significantly smaller, necessitating multiple epochs within the same token budget, resulting in stagnated model improvements.
\end{itemize}

\subsection{Training and Evaluating DeepSeekMath-Base 7B}

This section details the development of DeepSeekMath-Base 7B, a foundational model characterized by advanced mathematical reasoning capabilities. The model leverages DeepSeek-Coder-Base-v1.5 7B \citep{deepseek-coder} as its initialization point and is trained on a dataset totaling 500B tokens. The composition of the training corpus is as follows: 56\% DeepSeekMath Corpus, 4\% AlgebraicStack, 10\% arXiv, 20\% Github code, with the remaining 10\% consisting of natural language samples in English and Chinese sourced from Common Crawl. With respect to training configurations, we follow the methodology outlined in Section \ref{sec:quality-policy}, but adjust the peak learning rate to 4.2e-4 and employ a batch size of 10M tokens.

We systematically evaluate the mathematical proficiency of DeepSeekMath-Base 7B. The assessment encompasses its capacity for generating complete mathematical solutions autonomously, its competence in utilizing tools to address mathematical tasks, and its ability to engage in formal theorem proving. Additionally, we present an overview of the model’s general aptitude in areas such as natural language understanding, logical reasoning, and programming.

\paragraph{Mathematical Problem Solving with Step-by-Step Reasoning}

DeepSeekMath-Base is assessed on mathematical problem-solving tasks utilizing few-shot chain-of-thought prompting \citep{cot}, spanning eight benchmarks across both English and Chinese. The evaluation set includes quantitative reasoning tasks—such as those found in GSM8K \citep{gsm8k}, MATH \citep{MATH}, and CMATH \citep{wei2023cmath}—as well as multiple-choice assessments like MMLU-STEM \citep{mmlu} and Gaokao-MathQA \citep{agieval}. The selected benchmarks reflect a wide spectrum of mathematical disciplines and span from basic to collegiate levels of difficulty.

As indicated in Table \ref{tab:base_math_cot}, among open-source base models, DeepSeekMath-Base 7B consistently achieves the highest scores across all eight evaluated benchmarks, surpassing widely-used models such as Mistral 7B \citep{mistral} and the more recently introduced Llemma 34B \citep{llemma}, the latter of which was specifically trained on the Proof-Pile-2 mathematics corpus \citep{llemma}. Remarkably, on the MATH dataset—recognized for its competitive rigor—DeepSeekMath-Base 7B exceeds other open-source base models by more than 10\% in absolute terms, and further outperforms Minerva 540B \citep{minerva}, a proprietary model that is 77 times larger, built upon PaLM \citep{palm} with additional fine-tuning on mathematical resources.

\paragraph{Mathematical Problem Solving with Tool Use} To assess models’ abilities in program-aided mathematical reasoning, we employ few-shot program-of-thought prompting \citep{pot,pal} on both GSM8K and MATH datasets. Here, models are directed to produce Python programs that can call upon libraries such as \textit{math} and \textit{sympy} to facilitate complex calculations. The correctness of the solution is determined by running the generated code and comparing its output to the expected answer. According to Table \ref{tab:base_math_pot_proof}, DeepSeekMath-Base 7B surpasses the prior top-performing open-source model, Llemma 34B, on these program-aided tasks.

\paragraph{Formal Mathematics}

Automating the process of formal proof writing not only guarantees mathematical rigor but also promotes efficiency—leading to increasing interest in recent years. We evaluate DeepSeekMath-Base 7B on the informal-to-formal proof translation task outlined in \citep{dsp_proof}, where the model receives an informal statement, its formal version, and an informal proof, and must generate a formal proof. The assessment is conducted on the miniF2F dataset \citep{minif2f}, focused on Olympiad-grade formal mathematics, by requiring the model to write formal Isabelle proofs for each instance via few-shot prompting. Consistent with the setup of \cite{dsp_proof}, we utilize the model to produce proof sketches, after which the Sledgehammer automated theorem prover \citep{sledgehammer} is used to fill any missing justifications. Table \ref{tab:base_math_pot_proof} reveals that DeepSeekMath-Base 7B attains a strong level of performance in proof autoformalization.

\paragraph{Natural Language Understanding, Reasoning, and Code}

We examine the model’s capabilities on natural language understanding using MMLU \citep{mmlu}, on complex reasoning using BBH \citep{bbh}, and on programming skills via HumanEval \citep{codex} and MBPP \citep{mbpp}. Table \ref{tab:understanding_reasoning_code} demonstrates that DeepSeekMath-Base 7B achieves considerable improvements over its predecessor, DeepSeek-Coder-Base-v1.5 \citep{deepseek-coder}, in both MMLU and BBH, which highlights the advantages conferred by intensive math-oriented training for general understanding and reasoning. Further, by incorporating code-related data throughout continued training, DeepSeekMath-Base 7B retains the coding proficiency of DeepSeek-Coder-Base-v1.5 on the programming tasks. Overall, DeepSeekMath-Base 7B displays a pronounced advantage over the general-purpose Mistral 7B \citep{mistral} across the trio of reasoning and coding evaluations.

\section{Supervised Fine-Tuning}

\subsection{SFT Data Curation}

We assemble an instruction-tuning dataset tailored for mathematics, featuring a range of English and Chinese questions that span multiple mathematical disciplines and complexity levels. Each question is accompanied by a solution presented in chain-of-thought (CoT) \citep{cot}, program-of-thought (PoT) \citep{pot,pal}, or tool-integrated reasoning format \citep{tora}. Our final curated dataset comprises 776K training instances.

\begin{itemize}[topsep=0pt]
    \item \textbf{English mathematical datasets}: We provide tool-integrated annotations for GSM8K and MATH datasets and further include a selected portion of MathInstruct \citep{MathInstruct}, along with the Lila-OOD \citep{lila} training data, which feature problems solved via CoT or PoT reasoning. This English collection encompasses a breadth of mathematical topics such as algebra, probability, number theory, calculus, and geometry.
    \item \textbf{Chinese mathematical datasets}: Our Chinese dataset consists of K-12 mathematics problems drawn from 76 unique subfields, for example, linear equations, each annotated with solutions employing both CoT and tool-integrated reasoning techniques.
\end{itemize}

\subsection{Training and Evaluating DeepSeekMath-Instruct 7B}

This section details the fine-tuning process of DeepSeekMath-Instruct 7B, which is derived from DeepSeekMath-Base through instruction tuning focused on mathematical reasoning. Training samples are randomly joined together to construct sequences up to a 4K token context length limit. We train the model over 500 iterations, utilizing a batch size of 256 and maintaining a constant learning rate of 5e-5.

The evaluation of mathematical competency is conducted both in tool-augmented and standalone scenarios, employing four quantitative reasoning benchmarks in both English and Chinese. We compare our model’s performance against state-of-the-art counterparts of its time:

\begin{itemize}[topsep=0pt]
    \item \textbf{Closed-source models} in our comparison include: (1) members of the GPT family, with GPT-4 \citep{gpt4} and GPT-4 Code Interpreter~\footnote{\url{https://openai.com/blog/chatgpt-plugins\#\#code-interpreter}} as the leading representatives, (2) Gemini Ultra and Pro \citep{gemini}, (3) Inflection-2 \citep{inflection-2}, (4) Grok-1~\footnote{\url{https://x.ai/model-card}}, as well as models recently introduced by Chinese companies such as (5) Baichuan-3~\footnote{\url{https://www.baichuan-ai.com}}, and (6) the most recent GLM-4~\footnote{\url{https://open.bigmodel.cn/dev/api\#glm-4}} from the GLM suite \citep{glm}.
\end{itemize}

These models are designed for broad applications, the majority of which have been subjected to various alignment techniques.

\item \textbf{Open-source models} comprise: foundational models such as (1) DeepSeek-LLM-Chat 67B \citep{deepseek-llm}, (2) Qwen 72B \citep{qwen}, (3) SeaLLM-v2 7B \citep{seallm}, and (4) ChatGLM3 6B \citep{chatglm3}, along with models enhanced for mathematical reasoning. This latter group includes (5) InternLM2-Math 20B~\footnote{\url{https://github.com/InternLM/InternLM-Math}}, an extension of InternLM2 specifically trained on mathematical tasks followed by instruction fine-tuning; (6) Math-Shepherd-Mistral 7B, which employs PPO training \citep{schulman2017proximal} on Mistral 7B \citep{mistral} combined with a process-supervised reward mechanism; (7) the WizardMath series \citep{wizardmath}, focused on advancing mathematical reasoning in Mistral 7B and Llama-2 70B \citep{llama2} using evolve-instruct—an instruction tuning method relying on AI-generated tasks—plus PPO training with training data predominantly drawn from GSM8K and MATH; (8) MetaMath 70B \citep{metamath}, which constitutes Llama-2 70B refined using an expanded version of the GSM8K and MATH datasets; (9) ToRA 34B \cite{tora}, developed by fine-tuning CodeLlama 34B for integrated mathematical tool use; and (10) MAmmoTH 70B \citep{MathInstruct}, featuring Llama-2 70B tuned with MathInstruct-focused instructions.

According to the data summarized in Table \ref{tab:sft_rl_math}, in the context where tool-based assistance is not permitted, DeepSeekMath-Instruct 7B demonstrates notable proficiency in step-by-step reasoning. Especially on the MATH benchmark targeting competition-level mathematics, this model exceeds the performance of all available open-source alternatives and outperforms most proprietary offerings (such as Inflection-2 and Gemini Pro) by at least 9\% in absolute terms. This lead holds even over models of much greater size (e.g., Qwen 72B) or those improved via specialized reinforcement learning approaches for mathematical reasoning (e.g., WizardMath-v1.1 7B). DeepSeekMath-Instruct delivers performance comparable to leading Chinese proprietary models such as GLM-4 and Baichuan-3 on MATH, but remains behind GPT-4 and Gemini Ultra.

In scenarios where models can exploit both natural language reasoning and programming-based tool utilization, DeepSeekMath-Instruct 7B achieves nearly 60\% accuracy on MATH, outpacing all other open-source systems. On additional benchmarks, it maintains competitiveness with DeepSeek-LLM-Chat 67B, a previously top-performing model of tenfold greater scale.

\section{Reinforcement Learning}

\subsection{Group Relative Policy Optimization} Recent work demonstrates that reinforcement learning (RL) is effective for further elevating mathematical reasoning skills in large language models post-Supervised Fine-Tuning (SFT) \citep{wang2023math,wizardmath}. We present Group Relative Policy Optimization (GRPO), an efficient and robust RL algorithm tailored for this task.

\subsubsection{From PPO to GRPO} Proximal Policy Optimization (PPO) \citep{schulman2017proximal} serves as a foundational actor-critic algorithm, frequently applied for RL-based fine-tuning of LLMs \citep{ouyang2022training}. Specifically, PPO seeks to maximize a clipped surrogate objective, expressed as:

\begin{equation}
\footnotesize
    \mathcal{J}_{PPO}(\theta) = \mathbb{E}{[q \sim P(Q), o \sim \pi_{\theta_{old}}(O|q)]} \frac{1}{|o|} \sum_{t=1}^{|o|} \min \left[ \frac{\pi_\theta(o_{t} | q, o_{<t})}{\pi_{\theta_{old}}(o_{t} | q, o_{<t})} A_{t}, \text{clip} \left( \frac{\pi_\theta(o_{t} | q, o_{<t})}{\pi_{\theta_{old}}(o_{t} | q, o_{<t})}, 1 - \varepsilon, 1 + \

Here, $\pi_{\theta}$ and $\pi_{\theta_{old}}$ denote the current and previous policy networks, respectively. The variables $q$ and $o$ represent questions and outputs, where the latter are sampled from the question dataset and the old policy $\pi_{\theta_{old}}$. The hyper-parameter $\varepsilon$ relates to the clipping mechanism employed in PPO to promote stable learning. The advantage term $A_t$ is estimated via Generalized Advantage Estimation (GAE) \citep{gae}, utilizing the reward sequence $\{r_{\ge t}\}$ alongside a trainable value function $V_{\psi}$. Consequently, PPO necessitates joint optimization of both a value function and a policy network, and to counteract the risk of overfitting to the reward model, a typical strategy involves appending a per-token KL regularization penalty from a reference model to the reward, as proposed in \citep{ouyang2022training}, formulated as

\begin{equation}
    r_{t} = r_\varphi(q, o_{\le t}) - \beta \log\frac{\pi_{\theta}(o_{t}|q, o_{<t})}{\pi_{ref}(o_{t}|q, o_{<t})},
\label{eq:PPO-reward}
\end{equation}

where $r_\varphi$ denotes the reward model, $\pi_{ref}$ is the reference (usually the initial SFT) model, and $\beta$ controls the magnitude of the KL term.

Training a value function in PPO requires maintaining another network of similar size to the policy network, thereby increasing both memory consumption and computational load. Moreover, in the context of reinforcement learning, the value function is utilized as a baseline to reduce the variance in advantage computation. However, in the large language model (LLM) setting, it is common that only the reward for the final token is provided by the reward model, rendering it more challenging to learn a per-token accurate value function. To overcome these limitations, we introduce Group Relative Policy Optimization (GRPO), visualized in Figure \ref{fig:grpo}, which dispenses with the requirement for an explicit value function by instead taking the mean reward across a set of sampled outputs generated in response to the same prompt as the baseline. Specifically, for every question $q$, GRPO collects a batch of outputs $\{o_1, o_2, \cdots, o_G\}$ sampled from the previous policy $\pi_{\theta_{old}}$, and then updates the policy by maximizing

\begin{equation}
\footnotesize
\begin{split}
    \mathcal{J}_{GRPO}(\theta) &= \mathbb{E}{[q \sim P(Q), \{o_i\}_{i=1}^G \sim \pi_{\theta_{old}}(O|q)]}  \\
    & \frac{1}{G}\sum_{i=1}^G\frac{1}{|o_i|} \sum_{t=1}^{|o_i|} \left\{ \min \left[ \frac{\pi_\theta(o_{i,t} | q, o_{i,<t})}{\pi_{\theta_{old}}(o_{i,t} | q, o_{i,<t})} \hat{A}_{i,t}, \text{clip} \left( \frac{\pi_\theta(o_{i,t} | q, o_{i,<t})}{\pi_{\theta_{old}}(o_{i,t} | q, o_{i,<t})}, 1 - \varepsilon, 1 + \varepsilon \right)  \hat{A}_{i,t} \right] - \beta \mathbb{D}_{KL}\left[\pi_{\theta} || \pi_{ref}\right]\right\} ,
\end{split}
\label{eq:GRPO-obj}
\end{equation}

where $\varepsilon$ and $\beta$ are tunable hyper-parameters, and $\hat{A}_{i,t}$ is the advantage that is computed using only the relative rewards of outputs within each sampled group (described in detail in later sections). The use of group-relative advantage computation in GRPO aligns naturally with the structure of reward models, which are often trained on pairwise (or listwise) comparisons among outputs for a given prompt. Furthermore, GRPO imposes KL regularization directly through a term in the loss, contrasting with PPO’s practice of adding it into the reward. This avoids increasing the complexity of calculating $\hat{A}_{i,t}$. Differing from the KL term in (\ref{eq:PPO-reward}), GRPO estimates KL divergence using the following unbiased estimator from \citep{

\begin{algorithm}[t]
  \small
  \caption{Iterative Group Relative Policy Optimization}
  \textbf{Input} Initial policy model $\pi_{\theta_{\text{init}}}$; reward models $r_\varphi$; set of task prompts $\mathcal{D}$;
  hyperparameters $\varepsilon$, $\beta$, $\mu$
  \begin{algorithmic}[1]
    \State Initialize policy model $\pi_\theta \leftarrow \pi_{\theta_{\text{init}}}$
    \For{iteration = 1, \dots, I}
       \State Set reference model $\pi_{ref} \leftarrow \pi_{\theta}$
      \For{step = 1, \dots, M}
      \State Draw a minibatch $\mathcal{D}_b$ from $\mathcal{D}$
      \State Copy current policy $\pi_{\theta_{old}} \leftarrow \pi_{\theta}$ 
      \State For every question $q \in \mathcal{D}_b$, generate $G$ outputs $\{o_i\}_{i=1}^G \sim \pi_{\theta_{old}} (\cdot \mid q)$
      \State Evaluate each $o_i$ with $r_\varphi$ to produce rewards $\{r_i\}_{i=1}^G$
      \State Use group-relative estimation to determine $\hat{A}_{i,t}$ for the $t$-th token in $o_i$
      \For{GRPO iteration = 1, \dots, $\mu$}
        \State Adjust policy parameters $\pi_{\theta}$ to maximize the GRPO objective (Equation \ref{eq:GC-GRPO})
      \EndFor
    \EndFor \State Continuously update $r_\varphi$ with replay-based learning. 
    \EndFor 
  \end{algorithmic}
  \textbf{Output} $\pi_\theta$
  \label{alg:iter-grpo}
\end{algorithm}

\subsubsection{Outcome Supervision RL with GRPO} To formally describe the process, for any given question $q$, a collection of responses $\{o_1, o_2, \cdots, o_G\}$ is drawn from the previous policy model $\pi_{\theta_{old}}$. Each response is evaluated by the reward model, which assigns scores $\mathbf{r}=\{r_1, r_2, \cdots, r_G\}$ accordingly. These scores are standardized by removing the mean of the group and dividing by their standard deviation. Under outcome supervision, the standardized reward is delivered only at the conclusion of each response $o_i$, and this value is assigned as the advantage $\hat{A}_{i,t}$ for every token in the sequence: $\hat{A}_{i, t} = \widetilde{r}_i = \frac{r_i - {\rm mean}(\mathbf{r})}{{\rm std}(\mathbf{r})}$. Policy optimization then proceeds by maximizing the criterion in equation (\ref{eq:GRPO-obj}).

\subsubsection{Process Supervision RL with GRPO} Since outcome supervision supplies only terminal rewards, it might fall short in guiding the learning process for more sophisticated mathematical reasoning tasks. Inspired by \cite{wang2023math}, process supervision is also employed, where rewards are granted at the end of every individual reasoning step. Specifically, for each prompt $q$ and generated sample set $\{o_1, o_2, \cdots, o_G\}$, a process reward model is used to assign stepwise rewards, yielding: $\mathbf{R} = \{ \{r_1^{index(1)},\cdots,r_1^{index(K_1)}\}, \cdots,  \{r_G^{index(1)},\cdots,r_G^{index(K_G)}\} \}$, where $index(j)$ marks the concluding token of the $j$-th reasoning step in each $o_i$, and $K_i$ is the number of reasoning steps in output $o_i$. As before, rewards are normalized using their mean and standard deviation: $\widetilde{r}_i^{index(j)} = \frac{r_i^{index(j)} - {\rm mean}(\mathbf{R})}{{\rm std}(\mathbf{R})}$.
Advantage estimates for each token are computed as the cumulative normalized rewards over all subsequent steps: $\hat{A}_{i, t} = \sum_{index(j) \ge t} \widetilde{r}_i^{index(j)}$.
The optimization then seeks to maximize the objective expressed in equation (\ref{eq:GRPO-obj

\subsubsection{Iterative RL with GRPO} During reinforcement learning, as training advances, the previous reward model may become inadequate for effectively supervising the evolving policy model. To address this, we experiment with iterative RL using GRPO. As detailed in Algorithm \ref{alg:iter-grpo}, the iterative GRPO approach involves regenerating the reward model's training data using samples produced by the policy model, and continuously updating the reward model through a replay strategy that retains 10\% of prior data. Subsequently, the reference model is updated to match the latest policy model, and the policy model itself undergoes further training with the updated reward model.

\subsection{Training and Evaluating DeepSeekMath-RL}

We apply reinforcement learning on DeepSeekMath-Instruct 7B, utilizing RL training data composed of chain-of-thought-format items drawn from GSM8K and MATH within the SFT corpus, totaling approximately 144K questions. In this phase, we omit other SFT-sourced questions to better analyze RL's effect on benchmarks that lack corresponding data during RL. Construction of the reward model training set is aligned with \citep{wang2023math}. The initial reward model is trained on DeepSeekMath-Base 7B, employing a learning rate of 2e-5. For GRPO training, the policy model uses a learning rate of 1e-6. The KL penalty coefficient is set to 0.04. For each item, we generate $64$ responses. The maximum sequence length is fixed at 1024, and the batch size is set to 1024. After each exploration round, the policy model receives just one update. Evaluation of DeepSeekMath-RL 7B follows the same protocols as DeepSeekMath-Instruct 7B. Within this setup, GSM8K and MATH tasks involving chain-of-thought reasoning are considered in-domain, whereas all remaining benchmarks are treated as out-of-domain.

Table \ref{tab:sft_rl_math} presents a comparative analysis of open- and closed-source models applying both chain-of-thought and tool-based reasoning strategies to benchmarks in English and Chinese. Key observations include: 1) Utilizing chain-of-thought reasoning, DeepSeekMath-RL 7B achieves accuracies of 88.2\% on GSM8K and 51.7\% on MATH, outperforming every open-source competitor in the 7B--70B parameter range and exceeding the performance of most closed-source alternatives. 2) Notably, DeepSeekMath-RL 7B’s training is confined to chain-of-thought instruction-tuning data derived from GSM8K and MATH, building upon DeepSeekMath-Instruct 7B. In spite of the limited scope of training data, it consistently surpasses DeepSeekMath-Instruct 7B in all metrics, highlighting the robust impact of reinforcement learning.

\section{Discussion}
Here, we summarize insights drawn from our experiences during both pre-training and reinforcement learning phases.

\subsection{Lessons Learnt in Pre-Training}

We begin with reflections on pre-training practices. Unless explicitly stated, the following adopts the configuration set forth in Section \ref{sec:quality-policy}. Importantly, references to the DeepSeekMath Corpus in this context pertain to the 89B-token dataset amassed during the project’s second data-gathering cycle.

\subsubsection{Code Training Benefits Mathematical Reasoning}

A commonly held, though largely untested, belief posits that code-oriented training can enhance reasoning capabilities. Our investigation provides supporting evidence within the context of mathematics: training with code data markedly improves a model’s proficiency in mathematical reasoning, irrespective of tool involvement.

To probe the influence of code training on mathematical reasoning abilities, we devised experiments encompassing both two-stage and single-stage training protocols:

\textbf{Two-Stage Training}

\begin{itemize}[topsep=0pt]
    \item \textbf{Code Training for 400B Tokens $\rightarrow$ Math Training for 150B Tokens}: DeepSeek-LLM 1.3B is first pretrained on 400B code tokens, subsequently followed by an additional 150B tokens focusing on mathematics.
    \item \textbf{General Training for 400B Tokens $\rightarrow$ Math Training for 150B Tokens}: As a baseline comparison, we replace code tokens with general domain tokens (sourced from a comprehensive general-purpose corpus by DeepSeek-AI) in the initial phase. This aims to clarify whether pretraining on code particularly improves later mathematical reasoning when compared to pretraining on broader text.
\end{itemize}

\textbf{One-Stage Training}

\begin{itemize}[topsep=0pt]
    \item \textbf{Math Training for 150B Tokens}: DeepSeek-LLM 1.3B is trained exclusively on 150B mathematics tokens.
    \item \textbf{Training on a mixture of 400B Code Tokens and 150B Math Tokens}: Given that math training directly following code training negatively impacts the model’s coding capabilities, we further examine whether a joint one-stage mixture of code and math tokens both enhances mathematical reasoning and alleviates catastrophic forgetting.
\end{itemize}

\paragraph{Results}
The downstream evaluation, as summarized in Table \ref{tab:code-math} and Table \ref{tab:code-math-general-eval}, provides insight into performance across these distinct training approaches.

Pretraining with code data is advantageous for program-assisted mathematical reasoning in both sequential (two-stage) and joint (one-stage) paradigms. According to Table \ref{tab:code-math}, initializing with code-focused training notably boosts the model’s proficiency on tasks such as GSM8K and MATH when tool use is allowed (i.e., via Python execution). Subsequent math-centered finetuning brings additional gains. Furthermore, when code and math tokens are combined and fed to the model in a single unified stage, there is a marked reduction in catastrophic forgetting (a challenge observed in two-stage regimes), while also maintaining strengths in both coding (see Table \ref{tab:code-math-general-eval}) and program-augmented math reasoning (Table \ref{tab:code-math}).

Code pretraining similarly benefits pure mathematical reasoning without tool reliance, though the improvements are more modest in the two-stage scenario. Early code-focused training streamlines the efficiency and overall outcome of later mathematical training, culminating in peak performance. Conversely, integrating both code and math tokens in a one-stage setting appears to negatively affect mathematical reasoning when tools are not involved. One possible explanation is that DeepSeek-LLM 1.3B, due to its constrained size, lacks the required capacity to optimally internalize and balance both modalities simultaneously.

\subsubsection{ArXiv Papers Exhibit Limited Benefit for Mathematical Reasoning Enhancement}

In many contemporary math-focused pre-training datasets, arXiv papers constitute a substantial portion \citep{minerva,gpt-f,llemma,mathpile}. Nonetheless, little detailed investigation has been dedicated to understanding their concrete influence on mathematical reasoning capabilities. Contrary to perhaps common intuition, our experimental outcomes suggest that arXiv papers offer minimal advantage for strengthening mathematical reasoning. To examine this, we evaluated models across different parameter scales—namely DeepSeek-LLM 1.3B and DeepSeek-Coder-Base-v1.5 7B \citep{deepseek-coder}—using two distinct arXiv datasets subjected to alternative preprocessing approaches:

\begin{itemize}[topsep=0pt]
    \item \textbf{MathPile} \citep{mathpile}: an 8.9B-token collection assembled through heuristic filtering and data cleansing, with scientific arXiv papers comprising over 85\% of the corpus;
    \item \textbf{ArXiv-RedPajama} \citep{redpajama}: a dataset of 28.0B tokens consisting of the full arXiv LaTeX sources, from which preambles, comments, macros, and references have been excised.
\end{itemize}

For the empirical studies, DeepSeek-LLM 1.3B underwent 150B tokens of training, while DeepSeek-Coder-Base-v1.5 7B was trained for 40B tokens using each of the aforementioned arXiv datasets. The results indicate that exclusive training on arXiv corpora leads to little to no improvement—and at times, even declines—in mathematical reasoning across a spectrum of mathematical evaluation benchmarks of varying levels of difficulty. These include quantitative problem sets like GSM8K and MATH (Table \ref{tab:arxiv-cot}), multiple-choice assessments such as MMLU-STEM (Table \ref{tab:arxiv-cot}), and formal mathematics tasks, exemplified by miniF2F (Table \ref{tab:arxiv_atp}).

Despite these observations, certain caveats must be noted, and conclusions should be regarded as preliminary. The scope of our analysis does not encompass:

\begin{itemize}[topsep=0pt]
    \item Potential effects of arXiv-derived tokens on specialized mathematical tasks not examined in this work, such as the informalization process where formal theorems or proofs are translated into informal language;
    \item Interactions or synergistic effects arising when arXiv content is combined with other data sources;
    \item Possible advantages that may emerge only when scaling models to much larger sizes.
\end{itemize}

Consequently, additional investigations will be necessary, and we reserve further study of these open questions for future work.

\subsection{Insights of Reinforcement Learning}
\subsubsection{Toward a Unified Framework} In the following, we propose a unified framework to compare and analyze several training algorithms, such as SFT, RFT, DPO, PPO, GRPO, and conduct empirical studies to probe the dynamics governed by this unified view. Typically, for any training approach, the parameter gradient with respect to $\theta$ may be represented as:

\begin{equation}
    \nabla_{\theta}\mathcal{J}_{\textcolor{red}{\mathcal{A}}}(\theta) = \mathbb{E}[\underbrace{(q,o) \sim \textcolor{red}{\mathcal{D}}}_{Data \ Source}]\left( \frac{1}{|o|} \sum_{t=1}^{|o|}  \underbrace{GC_{{\mathcal{A}}}(q, o, t, \textcolor{red}{\pi_{{rf}}})}_{Gradient \ Coefficient}  \nabla_{\theta}\log \pi_{\theta}(o_t | q, o_{<t})\right).
\label{eq:objective}
\end{equation}

This formulation incorporates three principal elements: 1) \textit{Data Source $\mathcal{D}$}, which specifies the dataset leveraged during model training; 2) \textit{Reward Function $\pi_{{rf}}$}, from which reward signals are derived; and 3) \textit{Algorithm $\mathcal{A}$}, responsible for integrating training data and reward signals into a gradient coefficient $GC$ that influences the extent of punishment or reward attributed to each data instance. Our discussion evaluates several prominent approaches within this consolidated framework:

\begin{itemize}
    \item  \textbf{Supervised Fine-tuning (SFT)}: SFT involves additional training of a pretrained model on datasets curated by humans.
    \item \textbf{Rejection Sampling Fine-tuning (RFT)}: In RFT, the SFT model undergoes further tuning with a subset of its own generated outputs (in response to SFT prompts), selecting those outputs that are verified as correct.
    \item \textbf{Direct Preference Optimization (DPO)}: DPO refines the SFT model by exposing it to alternative candidate outputs sampled from itself, optimizing with a pairwise preference-based DPO loss.
    \item \textbf{Online Rejection Sampling Fine-tuning (Online RFT)}: Unlike conventional RFT, Online RFT starts with the SFT-initialized policy, and iteratively fine-tunes using outputs sampled from the latest policy in an online setting.
    \item \textbf{PPO/GRPO}: Here, the policy is initialized with the SFT model and then continually enhanced by sampling and reinforcing outputs from the evolving policy in real time.
\end{itemize}

The elements of these approaches are outlined in Table \ref{tab:data_policy}. For a comprehensive derivation, see Appendix \ref{app:analysis-rl}.

\paragraph{Insight on Data Source} We classify data sources into two primary types: online sampling and offline sampling. Online sampling indicates that training samples are generated dynamically from the active exploration of the model being trained, whereas offline sampling utilizes data obtained from the static outputs of the initial SFT model. RFT and DPO adopt the offline approach, whereas Online RFT and GRPO utilize an online approach.

Referencing Figure \ref{fig:rl-analysis}, we observe that Online RFT delivers marked improvements over RFT on both evaluation benchmarks. In particular, Online RFT matches the performance of RFT during the early phases of training, but then substantially surpasses it as training progresses, highlighting the benefits of online training. This result aligns with intuition: at the outset, both the actor and the SFT model are nearly indistinguishable, and thus their sampled data differ minimally. However, as training advances, the actor’s samples deviate increasingly from the SFT model, amplifying the benefits of using up-to-date, model-specific data.

\paragraph{Insight on Gradient Coefficient} The transformation of the reward signal into a gradient coefficient serves to direct parameter updates within the model. In our experimental setup, reward functions are grouped as either `Rule’ or `Model.’ The `Rule’ category bases evaluation on whether the generated response is correct, while the `Model’ category employs a separately trained reward model (itself trained using the rule-based judgments) to assign scores to each response. Equations \ref{eq:GC-RFT} and \ref{eq:GC-GRPO} expose a fundamental distinction: GRPO uniquely modifies its gradient coefficient in proportion to the reward signal assigned by the reward model. This enables nuanced reinforcement or suppression depending on response quality. Online RFT, conversely, does not incorporate this adaptation; it treats all correct responses equivalently and refrains from penalizing incorrect answers, resulting in uniform positive reinforcement for correct outputs.

As illustrated in Figure \ref{fig:rl-analysis}, GRPO outperforms online RFT, underlining the effectiveness of modifying the positive and negative gradient coefficients. Moreover, GRPO+PS delivers better results than GRPO+OS, showcasing the advantage of incorporating granular, step-sensitive gradient adjustments. Our analysis also considers iterative reinforcement learning; specifically, we perform two rounds of iteration. As depicted in Figure \ref{fig:iter-rl}, these additional RL iterations yield substantial gains, with the most marked improvement observed during the initial iteration.

\subsubsection{Why RL Works?} In this study, reinforcement learning is applied using a subset of instruction tuning data, leading to considerable improvements over models solely trained with instruction tuning. To clarify the source of RL’s effectiveness, we assess both Pass@K and Maj@K accuracies of the Instruct and RL-augmented models across two benchmarks. The results, presented in Figure \ref{fig:maj-pass}, reveal that RL raises Maj@K performance without affecting Pass@K. This pattern suggests that RL contributes primarily by making the output distribution more resilient—\textbf{that is, the observed gains are likely a result of elevating correct responses within the TopK outputs, rather than a direct enhancement of the model’s core reasoning abilities.} Similarly, \citep{wang2023making} describes a \textbf{misalignment problem} in the reasoning capabilities of SFT models, finding that their performance on reasoning tasks can be improved via various preference alignment techniques \citep{yuan2023rrhf,song2023preference,wang2023making}.

\subsubsection{How to Achieve More Effective RL?}
\label{sec:effec_rl}
Our findings confirm that RL yields strong results for mathematical reasoning challenges, and we introduce a comprehensive paradigm for interpreting a range of representative training approaches. According to this framework, these methods can all be viewed as different manifestations of direct or approximated RL approaches. As formalized in Equation \ref{eq:objective}, the process is governed by three central components: Data Source, Algorithm, and Reward Function. We identify several avenues for advancing each of these components in future research.

\paragraph{Data Source} The data source constitutes the foundational input for every training approach. Specifically, within RL, it encompasses unlabeled questions paired with outputs drawn from the policy model. This work limits the data to questions used in the instruction tuning phase, generating outputs via basic nucleus sampling. We hypothesize this constraint may partially account for RL’s exclusive improvements to Maj@K performance. Moving forward, we plan to extend our RL procedures to cover out-of-distribution prompts and experiment with \textbf{more advanced sampling (decoding) mechanisms}, such as tree-search-based approaches \citep{yao2023tree}. Furthermore, \textbf{efficient inference strategies} \citep{xia-etal-2023-speculative,leviathan2023fast,kwon2023efficient,xia2024unlocking}, which significantly impact the exploration capability of policy models, are also expected to play a critical role.

\paragraph{Algorithms} Algorithms utilize data and the associated reward signal to compute the gradient coefficient necessary for adjusting model parameters. As indicated by Equation \ref{eq:objective}, virtually all contemporary approaches place complete \textbf{TRUST} in the reward function to modulate the conditional probability of generating specific tokens. Nonetheless, it is unrealistic to assume the infallibility of the reward signal, particularly in tasks characterized by high complexity. For instance, despite the rigorous annotation provided by expert annotators for the PRM800K datasets \citep{lightman2023let}, an estimated 20\% of annotations remain incorrect\footnote{\url{https://github.com/openai/prm800k/issues/12\#issuecomment-1728491852}}. Addressing this issue, we seek to investigate reinforcement learning algorithms capable of maintaining robustness in the presence of noisy reward signals. We anticipate that \textbf{WEAK-TO-STRONG} alignment methods \citep{burns2023weak} have the potential to fundamentally transform existing learning algorithms.

\paragraph{Reward Function} The reward function serves as the principal source of training feedback. Within the context of RL, the reward function typically takes the form of a neural reward model. We propose three key avenues for advancing reward models: 1) \textbf{Improving the generalization capacity of reward models.} Effective reward models should be capable of robustly generalizing to out-of-distribution prompts and sophisticated outputs; failing this, reinforcement learning may only serve to maintain the LLM distribution rather than augmenting its core performance; 2) \textbf{Capturing the reward model's uncertainty.} The modeling of uncertainty could facilitate connections between weak reward models and weak-to-strong alignment approaches; 3) \textbf{Constructing efficient, high-quality process reward models} capable of delivering granular feedback throughout complex reasoning tasks \citep{lightman2023let,wang2023math}.

\section{Conclusion, Limitation, and Future Work}
In this work, we introduce DeepSeekMath, a model that surpasses all existing open-source systems on the competition-level MATH benchmark, nearing the proficiency of proprietary solutions. Built upon DeepSeek-Coder-v1.5 7B, DeepSeekMath was continually trained for 500B tokens, with a considerable portion—120B math tokens—curated from Common Crawl data. Our comprehensive ablation analysis indicates that web-sourced content provides valuable, high-quality mathematical data, whereas arXiv does not offer the anticipated benefits. We propose Group Relative Policy Optimization (GRPO), an adaptation of Proximal Policy Optimization (PPO), which enhances mathematical reasoning with reduced memory requirements. Experimental findings demonstrate the efficacy of GRPO, even when DeepSeekMath-Instruct 7B achieves high scores on various benchmarks. Furthermore, we establish a unified perspective to interpret a range of existing techniques, and outline several promising avenues for optimizing reinforcement learning strategies.

While DeepSeekMath attains strong results on quantitative reasoning tasks, its performance in geometry and theorem-proving remains inferior compared to leading closed-source models. For example, during preliminary testing, the model struggled with problems involving geometric figures such as triangles and ellipses, possibly reflecting selection bias in both pre-training and fine-tuning data. Moreover, due to its smaller size, DeepSeekMath is less effective than GPT-4 in few-shot learning; GPT-4 benefits from few-shot prompts, whereas DeepSeekMath exhibits comparable results under zero-shot and few-shot conditions. Going forward, we aim to enhance our data selection methods to assemble higher-quality pre-training corpora, and intend to further explore the strategies described in Section \ref{sec:effec_rl} for advancing reinforcement learning techniques for LLMs.

\end{document}